\section{Mathematical Frameworks for Discovery}
\label{sec:math_framework}

The current implementation of $E^3$ focuses on continuous-time deterministic dynamical systems represented as JAX functions and evaluated with Diffrax ODE solves. Many real systems are stochastic or event driven, and extending the framework to those classes remains a central next step.

\subsection{Ordinary Differential Equations (ODEs)}
\subsection{Model Language and Hypothesis Space}
To prevent the generation of pathological or unidentifiable mathematical structures, $E^3$ is constrained to a strictly defined vocabulary of primitives. The allowable mathematical operators include basic arithmetic ($+, -, \times, \div$), polynomials up to degree 3, trigonometric functions ($\sin, \cos$), exponential functions ($\exp$), and bounded regime switches ($\text{sigmoid}$, $\text{step}$). 

Furthermore, the engine enforces strict physical constraints directly within the JAX/Diffrax compilation step:
\begin{enumerate}
    \item \textbf{Non-Negativity:} Parameters representing physical mass, concentrations, or rates are strictly bounded using \texttt{jnp.clip} or exponential mappings ($e^\theta$) to ensure they never fall below zero.
    \item \textbf{Conservation Laws:} The LLM is prompted to explicitly output derivative structures that sum to zero ($\sum \frac{dy_i}{dt} = 0$) when modeling closed compartmental systems.
    \item \textbf{Stiffness Mitigation:} To prevent the Tsit5 solver from hanging on divergent asymptotic limits, the integration horizon is strictly bounded, and \texttt{max\_steps} is constrained to 4096. Any proposal that triggers a \texttt{max\_steps} exception is immediately rejected.
\end{enumerate}

The backbone of $E^3$ relies on continuous-time deterministic evolution. For macroscopic physical systems, chemical kinetics, and large-scale population dynamics, the LLM proposes an ODE state derivative function:
\begin{equation}
    \frac{d\mathbf{y}}{dt} = f(\mathbf{y}, t; \theta)
\end{equation}
where $\mathbf{y}$ is the state vector, $t$ is time, and $\theta$ is the parameter set. The function $f$ is generated as raw JAX/Python code by the LLM and solved using the Diffrax library via the Tsitouras 5/4 (Tsit5) adaptive solver. By generating raw code rather than manipulating abstract syntax trees, the LLM can express JAX primitives such as \texttt{jnp.where} for discontinuous boundary conditions or \texttt{jnp.sin} for external periodic forcing.

\subsection{Future Extension: Poisson Point Processes}
Several target domains, especially epidemiology, would be better modeled with discrete stochastic arrival processes rather than deterministic ODEs. A future version of $E^3$ should support \textbf{Poisson Point Process} simulators with jump events:
\begin{equation}
    P(N(t + \Delta t) - N(t) = k) = \frac{(\int_t^{t+\Delta t} \lambda(\tau, \mathbf{y}; \theta) d\tau)^k e^{-\int_t^{t+\Delta t} \lambda(\tau, \mathbf{y}; \theta) d\tau}}{k!}
\end{equation}
where the time-varying intensity $\lambda(\tau, \mathbf{y}; \theta)$ could be dynamically constructed by the LLM based on system states.

This is not implemented in the current codebase. The current ABC-SMC engine uses deterministic simulations and Euclidean distances over trajectory or summary-statistic features.

\subsection{Future Extension: Stochastic Differential Equations}
For macroeconomic systems and micro-scale cellular biology, deterministic representations may be inadequate. A future implementation could support Ito Stochastic Differential Equations:
\begin{equation}
    d\mathbf{y}_t = \mu(\mathbf{y}_t, t; \theta)dt + \sigma(\mathbf{y}_t, t; \theta)dW_t
\end{equation}
where $W_t$ is a multidimensional Wiener process. In this paradigm, the LLM would propose both the deterministic drift term $\mu$ and diffusion term $\sigma$. Diffrax provides tools for this path, but the present experiments do not yet use an integrated SDE discovery loop.
